% 第 2 章 执行摘要与核心观点
% AStock House View — 必须是独立观点，不能只复述卖方

\section{AStock House View：黄金拐点的四主线策略}

2026 年全球存储产业正处于过去 15 年以来最确定的一次「结构性上行周期」。与 2016-2018 年、2020-2021 年两轮由手机和 PC 拉动的经典周期不同，本轮上行的核心需求来源——AI 训练/推理集群——具有更高的单位算力存储强度、更长的建设周期 (3-5 年) 以及更高的客户集中度。同时，供给端在经历了 2022-2024 三年的深度减产 discipline 后，原厂 (三星/SK 海力士/美光) 的 capex 纪律性显著提升，不再出现 2018 年那种产能集中释放导致的价格崩塌。

\begin{houseviewbox}[AStock 独立观点 (House View)]
\vspace{0.3em}
\textbf{核心论点一：HBM 是第一增长极，但 A 股直接受益标的稀缺，市场对国产 HBM 的节奏严重误判。}
HBM 2026E 全球 TAM \$380--420 亿，CAGR 68\%，是半导体行业增速最快的细分赛道。但 HBM 市场高度集中：SK 海力士一家独大 (55--60\%)，三星和美光追赶 (合计 40--45\%)，三家原厂的良率、产能和客户绑定构筑了极高的护城河。\textbf{国产 HBM 的实质可用时点至少在 2028 年以后}，当前 A 股市场部分「国产 HBM 概念股」的炒作基于错误的预期管理和节奏误判，投资者需警惕此类标的的估值回撤风险。\par
\vspace{0.3em}
\textbf{核心论点二：设备材料确定性最高（设备国产化率 10\%\(\to\)25\%，约 18 个月翻倍），澜起 RCD 全球份额 ≈70\% 卡位单机必选。}
不论 HBM 竞争格局谁赢谁输，全球存储全产业链（含原厂+设备+材料）2026E Capex 约 \$1200 亿，其中三大原厂（Samsung/SK/Micron）合计约 \$850--910 亿，加上 YMTC/CXMT 国内原厂投入，设备和材料环节的订单确定性极高。北方华创 (平台化刻蚀/薄膜)、中微公司 (3D NAND 介质刻蚀 + TSV)、拓荆科技 (ALD/CVD) 三家在 YMTC/CXMT 的份额已从 2023 年的 <10\% 快速提升至 2026E 的 20-35\%，叠加 BIS 管制下的国产化刚性替代，确定性高于任何设计端标的。澜起科技则凭借 DDR5 RCD 全球 约70\% 份额 + CXL 内存扩展控制器卡位，成为 A 股唯一直接受益于 AI 服务器单机内存量价齐升的标的。\par
\vspace{0.3em}
\textbf{核心论点三：2026Q3 起 DRAM 合约价进入加速上行通道，利基存储 (兆易/江波龙) 弹性最大。}
三大原厂 2026 年 capex 虽同比增长 10-30\%，但增量主要投向 HBM 和先进制程，标准 DDR5 和 DDR4 的 bit growth 被刻意压制。叠加 AI 服务器单机 DDR5 容量从 1TB 升至 1.5-2TB、PC 和手机端 DDR5 渗透率从 50-60\% 升至 70-80\%，2026Q3 起标准 DRAM 供需缺口将扩大至 -8---10\%，合约价 QoQ +10--+12\%。利基 DRAM (海外大厂主动退出) 和 SLC NAND (2D NAND 产能收缩) 的涨价弹性将显著高于标准品。\par
\vspace{0.3em}
\textbf{核心论点四：CXL 商业化元年的预期差最大，澜起/江波龙/深科技具备 2027 年的盈利上修弹性。}
市场当前仅将 CXL 作为「远期题材」，但 CXL 2.0 控制器已由澜起量产、Intel Emerald Rapids 和 AMD Genoa 全面支持 CXL、Smart Global 的 CXL 内存模组已通过头部云厂商认证。预计 2026H2 CXL 在高端 AI 服务器中的渗透率将从当前 <3\% 快速升至 8-10\%，2027 年进一步升至 18-22\%。澜起的 CXL 相关收入在 2026 年约 10 亿元 (管理层指引)，2027 年有望翻倍至 20 亿+，存在 15-20\% 的盈利上修空间。
\end{houseviewbox}

\needspace{36\baselineskip}
\begin{exhibitbox}[图 2-1 \quad AStock House View 四核心论点结构 / Four Core Theses Structure]
\centering
\resizebox{0.88\exhibitwidth}{!}{%
\begin{tikzpicture}[
  thesis/.style={rectangle, rounded corners=4pt, minimum width=6.5cm, minimum height=1.4cm, align=center, draw=#1, fill=#1!10, text=deepnavy, font=\bfseries\small, inner sep=4pt},
  subp/.style={rectangle, rounded corners=2pt, minimum width=3.0cm, minimum height=1.1cm, align=center, draw=#1!70, fill=#1!5, font=\scriptsize, inner sep=4pt, text=deepnavy},
  arrow/.style={->, thick, >=Stealth, draw=steelgrey}
]
% Four theses arranged as 2x2
\node[thesis=riskred] (TH1) at (0, 3.2) {论点一：HBM 是第一增长极 · A 股直接受益稀缺};
\node[thesis=navy] (TH2) at (9.0, 3.2) {论点二：设备/材料是二阶 Beta · 澜起是一阶 Alpha};
\node[thesis=riskamber] (TH3) at (0, 0) {论点三：2026Q3 DRAM 合约价加速上行 · 利基弹性最大};
\node[thesis=nvgreen] (TH4) at (9.0, 0) {论点四：CXL 商业化元年 · 预期差最大};

% Sub points for each
\foreach \y in {2.1, 1.6, 1.1} \node[subp=riskred] at (0, \y-0.5) {\strut};
\node[subp=riskred, text width=3.0cm] at (0, 1.6) {国产 HBM ≥2028 · {\faExclamationTriangle} 市场误判节奏};
\node[subp=riskred, text width=3.0cm] at (0, 1.05) {SK 海力士 55\%+ · 三家原厂护城河极高};

\foreach \y in {2.1, 1.6, 1.1} \node[subp=navy] at (9.0, \y-0.5) {\strut};
\node[subp=navy, text width=3.0cm] at (9.0, 1.6) {北华/中微/拓荆 · 国产化率 10\%→25\%};
\node[subp=navy, text width=3.0cm] at (9.0, 1.05) {澜起 RCD 70\% + CXL · 必选税};

\foreach \y in {-1.1, -1.6, -2.1} \node[subp=riskamber] at (0, \y+0.3) {\strut};
\node[subp=riskamber, text width=3.0cm] at (0, -1.6) {标准品 bit growth ≤20\% · 供给刚性};
\node[subp=riskamber, text width=3.0cm] at (0, -2.15) {Q3 缺口-8~-10\% · QoQ +10\%+};

\foreach \y in {-1.1, -1.6, -2.1} \node[subp=nvgreen] at (9.0, \y+0.3) {\strut};
\node[subp=nvgreen, text width=3.0cm] at (9.0, -1.6) {CXL 2.0 量产 · 高端 AI 服 8-10\%};
\node[subp=nvgreen, text width=3.0cm] at (9.0, -2.15) {澜起 CXL 收入 26 10亿 · 27 翻倍};

% Arrows from thesis to conclusion center
\node[draw=deepnavy, fill=deepnavy, text=white, rounded corners=4pt, minimum width=9.5cm, minimum height=1.4cm, text width=14.4cm, align=center, font=\bfseries\normalsize, inner sep=5pt] (CONCL) at (4.5, -3.4) {AStock 结论：2026H2-2027 量价齐升双击 · 核心组合目标回报 +25-40\%};
\foreach \t/\a in {TH1/-30, TH2/-150, TH3/30, TH4/150} { \draw[arrow, bend right=\a] (\t.south) to (CONCL.north); }

% Top label
\node[fill=deepnavy, text=white, font=\bfseries\small, rounded corners=2pt, inner sep=3pt, anchor=west] at (-1.5, 4.8) {\faFlagCheckered\enspace House View · 四核心论点 → 组合结论};
\end{tikzpicture}
}
\end{exhibitbox}

\section{市场共识 vs AStock 独立观点}

\begin{exhibitbox}[表 2-1 \quad 市场共识 vs AStock 独立观点 / Consensus vs AStock House View]
\centering
\small
\begin{tabularx}{\exhibitboxwidth}{>{\raggedright}p{2.4cm} >{\raggedright}p{5.6cm} >{\raggedright\arraybackslash}X}
\toprule
\textbf{议题} & \textbf{卖方共识 (Consensus)} & \textbf{AStock 独立观点 (House View)} \\
\midrule
\rowcolor{riskred!5}
\textbf{HBM 国产节奏} & 2026-2027 年国内有望实现 HBM2E 小批量量产，2028 年 HBM3 & \textbf{{\faExclamationTriangle} 国产 HBM 实质可用 ≥2028 (HBM3 级)，A 股相关标的当前估值已提前透支 2 年预期。长鑫/长存未披露任何 HBM 进展，市场预期严重超前。} \\
\midrule
\rowcolor{nvgreen!5}
\textbf{设备 capex 弹性} & 受益存储扩产，设备厂收入 +30\% YoY & \textbf{实际弹性更大}：长存/长鑫国产化率从 10\%\(\to\)25\% 提供额外增量；BIS 管制下的「提前备货 + 战略库存」放大短期订单；北方华创等 2026 收入增速或超 40\%。 \\
\midrule
\rowcolor{riskamber!5}
\textbf{价格周期幅度} & 2026 DRAM 均价 +25--35\%，温和上涨 & \textbf{标准品涨幅或超预期}：原厂 bit growth 纪律严格，AI 服务器挤占标准品产能；2026Q3 缺口 -8---10\% 对应 QoQ +10\%+，全年涨幅可能超 45\%；但 2027H1 HBM 产能释放后存在回调风险。 \\
\midrule
\rowcolor{accentblue!5}
\textbf{CXL 渗透节奏} & 2026 年「题材年」，2028 年规模贡献 & \textbf{预期差所在}：CXL 2.0 控制器已量产、生态快速成熟；2026H2 高端 AI 服务器渗透率 8-10\%；澜起 CXL 收入 2026 年 约10 亿元、2027 年翻倍，非远期概念。 \\
\midrule
\textbf{长电/通富 HBM 封装} & 先进封装龙头直接受益 HBM 封装需求 & \textbf{需拆分验证}：长电/通富当前先进封装营收中 HBM 相关占比实际 <3\%；CoWoS 主产能仍集中在 TSMC；AI 封装收入实质贡献需 2027H2+ 方可见效。 \\
\midrule
\textbf{估值溢价合理性} & A 股存储 40x PE vs 海外 15-25x 偏高 & \textbf{溢价部分合理}：国产替代增速差 (35\% vs 15\%) + 国内政策红利 + 未上市资产稀缺性溢价。但纯概念标的 (无实质 AI 敞口) 溢价不可持续。 \\
\bottomrule
\end{tabularx}
\end{exhibitbox}
\sourcenote{卖方共识基于 14 份券商 PDF 原件提取 (详见第 9 章)；AStock 独立观点基于 Industry Landscape + 本报告全链路模型推断。}

\section{核心声明可信度分级}

遵循来源治理框架，将本报告所有高影响声明按证据强度分级，向读者透明披露判断依据。

\begin{exhibitbox}[表 2-2 \quad 核心声明可信度分级 / High-Impact Claim Confidence Ranking]
\centering
\small
\begin{tabularx}{\exhibitboxwidth}{>{\centering\arraybackslash}p{0.6cm} >{\raggedright\arraybackslash\hspace{0pt}}X >{\centering\arraybackslash}p{1.3cm} >{\centering\arraybackslash}p{0.9cm} >{\raggedright\arraybackslash\hspace{0pt}}p{3.2cm}}
\toprule
\textbf{\#} & \textbf{声明内容} & \textbf{等级} & \textbf{来源数} & \textbf{核心依据} \\
\midrule
SC-1 & HBM 2026E 全球 TAM \$380--420 亿，CAGR 68\% & A+ (极强) & 4+ & SK Hynix 管理层指引 + TrendForce + 三家卖方交叉 \\
SC-2 & 澜起 DDR5 RCD 全球 约70\% 份额 & A+ & 3+ & 澜起 IR + 东海/开源/国元 三家卖方确认 \\
SC-3 & 全球存储全产业链 2026E Capex 约 \$1200 亿；三大原厂（Samsung/SK/Micron）合计约 \$850--910 亿 & A (强) & 4+ & 三家原厂 Q4/2025-2026Q1 CC 官方披露 + 产业链口径（设备/材料端） \\
SC-4 & 2026Q3 DRAM 供需缺口 -8---10\%，合约价 QoQ +10\%+ & B+ (较强) & 2+ & TrendForce 历史 + AStock bottom-up 模型 \\
SC-5 & 北方华创在长存刻蚀份额 >20\% & A (强) & 2+ & 公司 IR + 中国国际招标网公开数据 \\
SC-6 & 长存 232L 量产、290L 研发中 & A (强) & 2+ & 长存官网 + 武汉政府公开信息 \\
SC-7 & 国产 HBM 最早 2028 实质可用 & B (中等) & 1+ & 制程代差倒推 + {\faExclamationTriangle} 行业访谈综合 \\
SC-8 & 澜起 CXL 3.0 获三星/AMD 积极反馈 & B+ & 2+ & 澜起 IR + 开源/东海卖方确认 \\
SC-9 & CXL 2026E 全球 TAM \$35--45 亿 & B & 1+ & TrendForce + AStock bottom-up 测算 \\
SC-10 & BIS 2026Q3 HBM 单独禁售概率 45--55\% & C (推断) & 1+ & Reuters 传闻 + AStock 地缘政治评估 \\
\bottomrule
\end{tabularx}
\end{exhibitbox}
\sourcenote{可信度等级：A+ (L1×3+交叉验证) / A (L1×2+ 或 L1+L3 强交叉) / B+ (L2+L3 双源) / B (单一 L3 或 L2 独立) / C (模型推断 + 单一弱源)。详见附录 A 表 A-1 来源注册。}

\vspace{0.4em}
\noindent\textbf{投资含义}：10 条核心声明中 7 条达 A/B+ 级 (证据扎实)，仅 SC-7 (国产 HBM 节奏)、SC-9 (CXL TAM 规模) 和 SC-10 (管制概率) 处于 B/C 级。A 类声明对应的投资机会 (澜起/设备三强/HBM TAM 规模) 是本报告高置信度的核心仓位；B/C 类声明对应的方向性判断 (国产 HBM 规避、CXL 超预期弹性) 则为组合提供 alpha 来源——若 CXL 声明兑现，澜起目标价存在额外 15-20\% 上修空间；若国产 HBM 声明被市场修正，相关标的存在 20-30\% 回撤风险。

从声明置信度与投资仓位的映射关系看，本报告遵循「高置信度→高仓位」原则：7 条 A/B+ 声明对应的投资主题 (HBM TAM 规模、澜起 RCD 份额、原厂 capex 总量、DDR5 涨价弹性、设备国产化率跃升) 合计对应约 75\% 的核心组合仓位，而 3 条 B/C 声明对应的主题 (CXL 弹性、国产 HBM 规避、BIS 管制判断) 仅对应约 25\% 的卫星/主题仓位和对冲权重。这一「置信度-仓位」映射框架确保了即使 B/C 级声明被后续证伪 (如 CXL 渗透低于预期、或 BIS 管制实际未落地)，组合整体回报仍受高置信度 A 类声明支撑。\par
具体到操作层面，建议投资者分三批建仓：\textbf{第一批 (当前，约 40\% 目标仓位)} 配置澜起 + 北方华创 + 中微 (A 类声明最高置信度 3 只)；\textbf{第二批 (2026 年 7 月中报季前后，约 35\%)} 配置 DDR5 涨价受益的江波龙/深科技、CXL 预期差的澜起加仓、设备弹性的拓荆；\textbf{第三批 (2026Q3 BIS 管制靴子落地后，约 25\%)} 根据管制实际力度在长电/通富 (先进封装方向) 和南大光电 (ArF 光刻胶方向) 之间选择。若 BIS 管制实际未升级，则第三批仓位可超配至 30\%，主要加仓设备国产化率跃升的弹性标的。\par
\textbf{声明置信度的季度再评估机制}：每季度末重新评估 10 条核心声明的证据等级，若 SC-1 (HBM TAM) 因 NVIDIA Rubin 出货低于预期从 A+ 降至 A，则 HBM 相关仓位下调 5 个百分点；若 SC-7 (国产 HBM 节奏) 因长存实际披露 HBM 样品从 B 升级为 B+，则国产 HBM 规避池的部分标的可移入观察池。所有升级/降级决策均需交叉验证至少 2 个独立数据源。
